\documentclass[../../main.tex]{subfiles}

\begin{document}

\subsection*{\texttt{edep\_reader}}
  \texttt{edep\_reader} è un dato Complex, Unique e Context, il cui compito
  è rappresentare le strutture dati contenute nel output di EDep-Sim.

  In quanto Complex è responsabile del proprio input; passando un file .root
  nel json di configurazione la sua \texttt{factory} lo apre e cerca un
  \texttt{TTree} chiamato ``EDepSimEvents'', se lo trova controlla che
  contenga un branch di tipo \texttt{TG4Event}, una struct \textit{data only}
  di EDep-Sim contenente le informazioni di un singolo spill.

  Per ogni \texttt{context\_id} \texttt{edep\_reader} espone due interfacce:
  la prima è un \texttt{const~TG4Event~*} per avere compatibilità
  con eventuale codice preesistente che accede direttamente alle strutture dati
  originali, la seconda invece è una struttura ad albero che riproduce la
  gerarchia dell'interazione.
  Questa nuova struttura dati originale di SANDReco è stata creata per
  facilitare le iterazioni sulle particelle dato che nel formato
  \texttt{TG4Event} sono salvate tutte in un unico \texttt{std::vector}
  e i loro rapporti di parentela sono indicati con gli indici delle loro
  posizioni nel vector stesso.
  Inoltre durante la costruzione della struttura ad albero, utilizzando il
  \texttt{root\_tgeomanager} per accedere alla geometria, vengono aggiunte
  informazioni sui detector attraversati dalle varie tracce.

\end{document}